\documentclass[11pt,a4paper]{article}

\usepackage[margin=2.5cm]{geometry}
\usepackage[T1]{fontenc}
\usepackage[utf8]{inputenc}
\usepackage{lmodern}
\usepackage{amsmath, amssymb}
\usepackage{booktabs}
\usepackage{graphicx}
\usepackage[hidelinks]{hyperref}
\usepackage{microtype}

% Long note fields in the bibliography would otherwise overrun the measure.
\setlength{\emergencystretch}{3em}

\usepackage[style=authoryear, backend=biber, maxbibnames=99, sorting=nyt]{biblatex}
\addbibresource{references.bib}

\title{Revised Research Proposal\\[0.5em]
  \large Label-Free Irreversibility as a Safety Signal in World-Model Latents\\[0.3em]
  \large The Safety Dial: Deployment-Time Conservatism}
\author{Sachin Mohan \\ Student No. 2699183 \\
  Supervised by Geraud Nangue Tasse \\ University of the Witwatersrand}
\date{September 2026}

\begin{document}

\maketitle

\begin{abstract}
Latent safety filters removed the hand-tuned penalty weight from safe reinforcement
learning, but they did not remove supervision: the failure set that seeds the reachability
computation is a classifier trained on human-annotated failure observations, so the filter
cannot see a failure mode nobody labelled. This proposal asks whether that failure set can
instead be derived from the dynamics, using \emph{irreversibility} as the criterion, and
whether the resulting conservatism can be exposed as a threshold an operator sets at
deployment rather than a constant an engineer fixes during training.

The central methodological change from the August 2026 proposal is the evaluation
environment. The Safety-Gymnasium locomotion suite carries two distinct unsafe events: a
velocity-limit violation, which is extrinsic, recoverable, and is the cost signal every
published safe-RL baseline optimises; and falling over, which is intrinsic, irreversible,
and is reported through episode termination rather than through the cost channel. These two
signals are produced by different mechanisms, which provides a labelled mode and an
unlabelled mode inside a single recognised benchmark, and therefore a direct test of whether
a label-free method detects an unsafe event that a supervised one structurally cannot. A
feasibility probe is reported. The work is scoped to an Honours thesis completing in
November 2026 and a workshop submission in February 2027.
\end{abstract}

\section{What changed, and why}

The August 2026 proposal committed to label-free latent reachability, with OGBench-Cube as
the primary environment and a bespoke annotation harness supplying ground-truth
recoverability labels for evaluation. That commitment is retained. Three things change.

\paragraph{Environment.} Safety-Gymnasium was previously rejected on the grounds that its
hazards are recoverable. That reasoning was correct for the navigation suite, where hazards
are regions marked on the floor with an assigned cost and nothing in the dynamics
distinguishes them. It does not hold for the locomotion suite, where the robot can fall and
cannot get up. Section~\ref{sec:observation} sets out the distinction and the evidence.
Adopting a recognised benchmark supplies published baselines and removes the need to
hand-build an annotation set, which was a substantial and bias-prone piece of work.

\paragraph{Scope.} The August proposal carried four hypotheses across five stages, and a
subsequent internal planning note carried four further hypotheses across six stages. Neither
set has been reduced by experiment, and the combination is not deliverable in the time
available. This document commits to three hypotheses and states the gate that kills each.

\paragraph{Venue.} The target is an ICLR 2027 workshop rather than the main track. The main
track abstract deadline of 18 September 2026 passed deliberately without a submission: the
direction was not settled and a rushed abstract was not worth the slot. Workshop papers are
non-archival and ICLR's dual-submission policy permits them, so this preserves a later full
submission. Section~\ref{sec:timeline} gives the schedule.

\section{Background and the gap}

Safe reinforcement learning has historically specified unsafe behaviour through a penalty
term in the reward or a cost function under a constraint. Both require a human to quantify
harm in advance, and the resulting policy is sensitive to that quantity. ROSARL
\parencite{nanguetasse2023rosarl} reduces the problem by deriving a sufficient penalty from
intrinsic quantities rather than tuning it, but the operating point remains a single scalar
fixed before deployment.

Latent safety filters \parencite{nakamura2025generalizing} take a different route. They
compute a Hamilton-Jacobi backward reachable tube \parencite{bansal2017hamilton} in the
latent space of a learned world model, and use it to override a task policy when failure
becomes inevitable. This removes the penalty weight. UNISafe \parencite{seo2025unisafe} adds
an uncertainty axis so that the filter distrusts its own rollouts near the boundary, and
AnySafe \parencite{agrawal2025anysafe} makes the constraint adjustable at runtime through
conformal similarity.

The gap is that every one of these seeds its reachability computation with a
\emph{classifier over labelled failure observations}. The filter is therefore blind to any
failure mode that was not annotated. Replacing the penalty weight with a labelled failure set
relocates the specification problem rather than solving it.

Irreversibility is the natural label-free replacement. A transition is unsafe when no
available action sequence returns the system to where it was. This is a property of the
transition structure, requires no annotation, and is computable only with a world model,
since it depends on rollouts of actions that were never taken. Reversibility-aware
reinforcement learning \parencite{grinsztajn2021turning} arrives at the same criterion from
the alignment side, using a precedence classifier over unlabelled trajectories.

A second gap concerns the substrate. Every latent safety filter published to date runs on a
\emph{reconstructive} world model, where the latent is trained to support pixel
reconstruction and therefore retains the information a safety signal needs more or less by
construction. A joint-embedding predictive architecture (JEPA) is trained to keep what is
predictable and discard the rest. Whether a non-reconstructive latent retains what safety
needs is an open question, and it is the representation half of this proposal.

\section{The key observation}
\label{sec:observation}

The Safety-Gymnasium locomotion suite \parencite{ji2023safety} defines cost as a velocity
threshold: cost accrues when the agent's forward velocity exceeds a fixed limit. Falling
over is not part of the cost signal at all. It is reported through the \texttt{terminated}
flag inherited from the underlying MuJoCo locomotion task, because a fallen Hopper or Walker
is in an unhealthy state from which the episode cannot continue.

The two events therefore differ on every axis that matters here:

\begin{table}[h]
\centering
\begin{tabular}{lll}
\toprule
 & \textbf{Mode A: velocity violation} & \textbf{Mode B: falling} \\
\midrule
Channel        & \texttt{cost}                & \texttt{terminated} \\
Origin         & extrinsic, hand-set limit    & intrinsic, physical \\
Reversible     & yes, decelerate              & no, the robot cannot rise \\
In cost signal & yes                          & no \\
Baselines optimise it & yes                   & no \\
\bottomrule
\end{tabular}
\caption{The two unsafe events in the Safety-Gymnasium locomotion suite.}
\end{table}

This yields, inside a single recognised benchmark, exactly the structure the zero-shot
transfer experiment requires: a labelled failure mode that the supervised baseline is trained
on, and an unlabelled, genuinely irreversible failure mode on which both methods are
evaluated. No bespoke annotation harness is needed, because the simulator reports both modes
directly and they can be reserved for evaluation.

\subsection{Feasibility probe}

A short probe was run to confirm the two channels are mechanically distinct. Twenty episodes
per condition, two policies: uniform random actions, and a crude directed policy applying
maximum forward torque as a proxy for a return-seeking agent.

\begin{table}[h]
\centering
\begin{tabular}{llccc}
\toprule
Environment & Policy & Episodes with fall & Episodes with cost & Mean length \\
\midrule
SafetyHopperVelocity-v1   & random       & 20/20 & 1/20  & 24.4 \\
SafetyHopperVelocity-v1   & full forward & 20/20 & 20/20 & 21.9 \\
SafetyWalker2dVelocity-v1 & random       & 20/20 & 0/20  & 21.5 \\
SafetyWalker2dVelocity-v1 & full forward & 20/20 & 0/20  & 207.1 \\
\bottomrule
\end{tabular}
\caption{Feasibility probe. Falls are frequent; cost is frequently absent.}
\end{table}

Two readings, and the second matters more than the first.

The structural point holds unconditionally: falling and cost are produced by different
mechanisms, so a filter seeded from the cost signal has no information about falls. Walker2d
under the directed policy is the clearest case, walking for an average of 207 steps, never
once triggering cost, and falling in every episode.

The empirical point is weaker and is stated plainly. Hopper under the directed policy
triggers cost in every episode as well as falling, so the two signals \emph{co-occur} in that
regime and the modes are not cleanly separable there. Disjointness is therefore a property of
the robot and the policy, not of the benchmark as a whole. Under a random or a crude policy
the robots also fall almost immediately, which is not the regime of interest.

\paragraph{Gate 0.} Before any comparison is run, fall rate and cost rate must be measured
under \emph{trained} policies (unconstrained PPO and a Lagrangian-constrained variant) on each
candidate robot, and the environment selected on measured disjointness. If no robot shows
falls at a usable rate under a trained policy while cost remains rare, the transfer
experiment is void and the environment choice must be revisited. This mirrors the
dataset-composition control in the August proposal, which exists for the same reason.

\section{Research question and hypotheses}

\paragraph{Research question.} Can an unsafe event be detected from irreversibility in the
dynamics of a learned world model, without failure labels, when that event is absent from the
benchmark's own cost signal and therefore from every supervised baseline trained on it; and
does a non-reconstructive latent retain the information required to do so?

\begin{description}
  \item[H1 (Representation).] Information sufficient to predict an imminent irreversible
    failure is linearly decodable from world-model latents, and degrades measurably with
    predictor rollout depth. Frozen general-purpose features retain more of it than an
    end-to-end predictive latent, because predictive training discards the fine pose and
    balance detail that signals an imminent fall.
  \item[H2 (Detection).] A label-free irreversibility score computed in that latent separates
    pre-failure states from matched recoverable states at AUROC $\geq 0.8$, using no failure
    labels at any point in its construction.
  \item[H3 (Transfer and dial).] A filter built on H2 reduces the irreversible-failure rate
    that the cost signal ignores, at a task-return cost selected at deployment by moving a
    single threshold, recovering a safety-performance frontier that penalty-based baselines
    require one training run per operating point to trace.
\end{description}

H1 is the representation claim, H2 the estimator, H3 the control claim. H1 and H3 are the
two the project is built around; H2 is the bridge between them.

\section{Method}

\paragraph{Stage 0: harness and environment selection.} Collect trajectories from the
candidate locomotion environments under random, trained-unconstrained and
trained-constrained policies, recording pixels, simulator state, action, cost, termination,
episode identifier and step index. Rendering is at $224 \times 224$ to match the existing
world-model preprocessing. Run Gate 0 and fix the environment.

\paragraph{Stage 1: representation probe (gating).} Train linear and small MLP probes from
latents to quantities that predict a fall (torso height, pitch angle, angular velocity), on
encoder latents and on predicted latents at rollout depths $1 \ldots H$. Report $R^2$ and the
degradation curve. Compare substrates. \textbf{Kill test:} if the fall-predictive quantity is
not decodable at rollout depth, no rollout-based irreversibility estimate in that latent can
work, and the project reports that as the finding and falls back to the probe subspace.

\paragraph{Stage 2: irreversibility estimation.} Two estimators developed in parallel. The
\emph{rollout} estimator samples $K$ action sequences from a fixed distribution, rolls them
forward through the predictor, and scores $\rho = \min_j d(\hat{z}_j, z_t)$; fixing the action
distribution is what makes differences reflect dynamics rather than remaining control
authority. The \emph{precedence} estimator \parencite{grinsztajn2021turning} trains a
classifier on unlabelled trajectories to predict which of two latents came first, and treats
confident asymmetry as a marker of irreversibility; it requires no rollout at inference.
Forward dispersion is implemented as the comparison expected to lose. Evaluation is by AUROC
against the simulator's termination flag, held out and never used in training.

Two controls are mandatory and are designed before the experiment is run. States with reduced
control authority but no damage must not score as irreversible. The rollout horizon $H$ must
be swept to establish that the score does not simply track predictor drift.

\paragraph{Stage 3: supervised baseline and transfer.} Train a latent safety filter in the
UNISafe style \parencite{seo2025unisafe} on Mode A labels only. Train the proposed method on
the identical trajectories with no labels. Evaluate both on Mode B. Report dataset
composition first: the shared unlabelled set must contain Mode B instances or the result is
void.

\paragraph{Stage 4: the dial.} Sweep the threshold at deployment on a single trained model.
Report irreversible-failure rate against task return as the threshold varies, with conformal
calibration supplying a coverage statement, and compare the compute required to trace the
whole frontier against penalty-swept baselines that need one training run per point.

\subsection{Substrates}

Three, which together form the representation axis of H1.

\begin{itemize}
  \item \textbf{Frozen general-purpose features with a small action-conditioned predictor}, in
    the DINO-WM style \parencite{zhou2024dinowm}. Only the predictor is trained, which is
    cheap and carries little risk, and the encoder discards nothing task-specific. This is
    the workhorse.
  \item \textbf{An end-to-end JEPA} trained on the selected environment, in the LeWM style
    \parencite{maes2026lewm}, whose isotropic latent follows from SIGReg
    \parencite{balestriero2025lejepa}. This is the substrate H1 predicts will lose.
  \item \textbf{Released LeWM checkpoints on Push-T}, retained as the sanity substrate where
    the existing pipeline is known to work.
\end{itemize}

World-model training is infrastructure, not contribution. The substrate comparison is
deliberately not load-bearing: if the end-to-end JEPA cannot be trained to a usable standard
in the time available, H1 is reported on the frozen-feature substrate alone and the
comparison is dropped.

\section{Evaluation and baselines}

Safety-Gymnasium supplies the baselines that make the result legible: PPO-Lagrangian, CPO
\parencite{achiam2017cpo}, TRPO-Lagrangian and related constrained methods, together with
SafeDreamer \parencite{huang2024safedreamer} as the model-based comparison. The supervised
latent safety filter of Stage 3 is the baseline the central claim is measured against.

Primary metrics: AUROC of the irreversibility score against held-out termination; Mode B
failure rate under each filter; task return; and compute to trace the safety-performance
frontier. The penalty-sweep frontier already measured in pilot work
(Section~\ref{sec:prelim}) is the comparison for the last of these.

\section{Preliminary evidence}
\label{sec:prelim}

Pilot work on Push-T with a released LeWM checkpoint (192-dimensional latents) predates this
revision but establishes two things the plan depends on.

\paragraph{Physical quantities are decodable from JEPA latents, including imagined ones.}
Probes from LeWM latents to pusher position over 100k frames reach $R^2 = 0.965$ (ridge, RMSE
18.9 px) and $R^2 = 0.999$ (MLP, RMSE 3.2 px). On \emph{predicted} latents at rollout depth
five the MLP probe still reaches $R^2 = 0.98$, with radial error of median 10.2 px, 95th
percentile 35.5 px and maximum 112 px. Physical state survives the predictor, which is what
Stage 1 generalises to fall-predictive quantities.

\paragraph{Penalty-based safety fails in exactly the way the dial is meant to address.} A
penalty-augmented CEM planner was swept over seven penalty weights and five seeds against a
calibrated hazard region. At zero penalty the planner violated on 27\% of steps and finished
16.4 px from the goal. At every non-zero penalty the final block error was 70.7 px, which is
precisely the start-to-goal distance, meaning the block never moved at all, and the pusher
left the arena in every run. Safety was bought by abandoning the task at every operating point
that was not zero. This is a first-hand instance of the fragility that motivates moving the
operating point out of training and into deployment.

\section{Risks and limitations}

\begin{itemize}
  \item \textbf{Mode disjointness may not survive trained policies.} This is the largest risk
    and Gate 0 exists to detect it early. The Hopper result above already shows the modes can
    co-occur under an aggressive policy.
  \item \textbf{Irreversibility is not harm.} Task progress can be irreversible. The scope is
    limited to environments whose harmful failures are the irreversible ones, and this is
    stated rather than defended.
  \item \textbf{Horizon-bounded.} What is measured is ``cannot return within $H$'', not
    ``cannot return ever''.
  \item \textbf{Predictor reliability near the boundary.} Rollouts close to failure are where
    the predictor is least trustworthy, which may require an uncertainty gate in the style of
    UNISafe.
  \item \textbf{Probe supervision.} Probes are supervised on generic physical quantities, never
    on the safety concept. This is stated explicitly wherever probe results are reported, and
    the probe-space variant is a fallback, not the headline.
  \item \textbf{Simulation only.} The irreversibility argument is strongest as a claim about
    physical systems, and no physical validation is in scope.
  \item \textbf{Preprint rate.} The area is moving quickly. Novelty is rechecked against arXiv
    and OpenReview before submission.
\end{itemize}

\section{Timeline and infrastructure}
\label{sec:timeline}

\begin{table}[h]
\centering
\begin{tabular}{ll}
\toprule
Period & Work \\
\midrule
Late September 2026 & Stage 0 harness, Gate 0, environment fixed \\
October 2026        & Stage 1 probes and gating decision; Stage 2 estimators begun \\
Early November 2026 & Stage 2 AUROC with both controls \\
November 2026       & Thesis write-up \\
December 2026       & Stage 3 supervised baseline and transfer \\
January 2027        & Stage 4 dial and frontier; results frozen mid-January \\
1 February 2027     & ICLR 2027 workshop submission \\
\bottomrule
\end{tabular}
\end{table}

The accepted ICLR 2027 workshop list is announced on 29 November 2026; individual calls
cluster in late January and early February. Workshops run 29 to 30 April 2027 in San
Francisco. Short-paper tracks of three to five pages are mandatory at every ICLR workshop,
with full tracks typically four to nine pages.

\paragraph{Infrastructure note.} Safety-Gymnasium 1.0.0 is pinned to Python $\leq 3.10$,
gymnasium 0.29 and numpy $< 2$, while the existing world-model stack requires Python 3.11,
gymnasium 1.3 and numpy 2.x. These cannot share an environment. The intended arrangement is
two environments with an offline handoff: trajectory collection in the Safety-Gymnasium
environment, written to HDF5; world-model training and evaluation in the main environment.
Headless rendering is confirmed working through EGL at the required resolution.

\printbibliography

\end{document}
